\chapter{Wortstellung}
\label{cha:Wortstellung}

Die Wortstellung im Deutschen folgt bestimmten Regeln, die über die Grundpositionen hinausgehen.

\section{Stellung des Subjekts}
\label{sec:SubjektStellung}

\begin{itemize}
    \item \textbf{Position 1}: Normal: Subjekt
    \item \textbf{Position 1}: Topikalisierung: Andere Elemente
        \begin{itemize}
            \item Heute \textbf{geht} er einkaufen.
            \item Im Winter \textbf{schneit} es oft.
            \item Mir \textbf{geht} es gut.
        \end{itemize}
\end{itemize}

\section{Stellung der Negation}
\label{sec:NegationStellung}

\begin{itemize}
    \item \textbf{nicht}: Vor dem Verb oder am Ende
        \begin{itemize}
            \item Er \textbf{kommt nicht}.
            \item Er \textbf{hat} nicht \textbf{kommen} können.
        \end{itemize}
    \item \textbf{kein}: Vor dem Nomen
        \begin{itemize}
            \item Ich habe \textbf{kein} Geld.
        \end{itemize}
\end{itemize}

\section{Stellung von Adverbien}
\label{sec:AdverbStellung}

\begin{enumerate}
    \item Temporal (Wann?): oft am Anfang oder vor dem Verb
        \begin{itemize}
            \item \textbf{Heute} komme ich. / Ich komme \textbf{heute}.
        \end{itemize}
    \item Lokal (Wo?): nach dem Verb
        \begin{itemize}
            \item Ich bin \textbf{in Berlin}.
        \end{itemize}
    \item Modal (Wie?): vor dem Verb
        \begin{itemize}
            \item Er spricht \textbf{schnell}.
        \end{itemize}
\end{enumerate}

\section{Inversion (Umstellung)}
\label{sec:Inversion}

\begin{itemize}
    \item \textbf{Nach „und", „oder", „aber"}: Subjekt folgt
        \begin{itemize}
            \item Und \textbf{er} kam nach Hause.
            \item Aber \textbf{ich} blieb.
        \end{itemize}
    \item \textbf{Nach „da" (weil)}: Subjekt folgt
        \begin{itemize}
            \item Da \textbf{es} regnete, blieben wir.
        \end{itemize}
    \item \textbf{Nach „nicht nur...sondern auch"}: Subjekt folgt
        \begin{itemize}
            \item Nicht nur \textbf{ich}, sondern auch \textbf{er} kam.
        \end{itemize}
\end{itemize}

\chapterend